\exercisesection{4}

\begin{enumerate}
\def\labelenumi{\arabic{enumi}.}
\tightlist
\item
  Let A be an integrally closed Noetherian domain, and V a vector space of finite rank over the field of fractions of A. Show that, if \((\mathbf{M}_{\lambda})\) is any family of reflexive lattices of V all containing the same lattice N, the lattice \(M = \bigcap_{\lambda} M_{\lambda}\) is reflexive (consider the dual lattices \(M_{\lambda}^{*}\) ).
\end{enumerate}

* 2. Let A be an integrally closed Noetherian domain and E, F two finitely generated A-modules. Suppose that E is torsion-free, that \(\mathrm{Hom}_{\mathrm{A}}(\mathrm{E},\mathrm{F})\) is reflexive and that \(\mathrm{Ext}_{\mathrm{A}}^{1}(\mathrm{E},\mathrm{F}) = 0\) . Show then that F is reflexive. (Prove first that F is torsion-free; if \(\mathrm{T} = \mathrm{F}^{**}/c_{\mathrm{F}}(\mathrm{F})\) , calculate \(\mathrm{Ass}(\mathrm{Hom}_{\mathrm{A}}(\mathrm{E},\mathrm{T}))\) in two ways using the exact sequence:

\[
0 \rightarrow \operatorname{Hom} (\mathrm{E}, \mathrm{F}) \rightarrow \operatorname{Hom} (\mathrm{E}, \mathrm{F} ^ {* *}) \rightarrow \operatorname{Hom} (\mathrm{E}, \mathrm{T}) \rightarrow 0
\]

and Chapter IV, § 1, no. 4, Proposition 10.) *

\begin{enumerate}
\def\labelenumi{\arabic{enumi}.}
\setcounter{enumi}{2}
\item
  Let A be an integrally closed Noetherian domain, K its field of fractions, M a reflexive lattice of a vector space of finite rank over K and L a free lattice of V containing M. Show that there exists a free lattice \(L_{1}\) of V such that \(M = L \cap L_{1}\) . (Consider the finite set I of prime ideals p of A of height 1 such that \(L_{p} \neq M_{p}\) and the principal ideal domain \(S^{-1}A\) , where \(S = \bigcap_{p \in I} (A - p)\) ; show that there exist a free lattice \(L_{0}\) of V such that \(M_{p} = (L_{0})_{p}\) for all \(p \in I\) and an \(s \in S\) such that \(M \subset s^{-1}L_{0} \subset L_{1}\) .)
\item
  Let k be a field and A the polynomial ring \(k[X, Y]\) .
\end{enumerate}

\begin{enumerate}
\def\labelenumi{(\alph{enumi})}
\item
  Let \((e_1, e_2)\) be the canonical basis of the A-module \(A^2\) and let E be the sub-A-module of \(A^2\) generated by \((X - Y)e_1\) , \(e_1 + Xe_2\) and \(e_1 + Ye_2\) ; let F be the monogenous submodule of E generated by \((X - Y)^2e_1\) . Show that the A-module \(M = E / F\) is not the direct sum of its torsion submodule and a torsion-free module.
\item
  Show that the torsion A-module A/AXY is not a direct sum of monogenous submodules of the form \(A/P_{i}^{n_{i}}\) , where the \(P_{i}\) are prime ideals of height 1.
\end{enumerate}

\begin{enumerate}
\def\labelenumi{\arabic{enumi}.}
\setcounter{enumi}{4}
\tightlist
\item
  Let A be an integrally closed Noetherian domain and M a finitely generated A-module. Show that, if a, b are two ideals of A, the A-module
\end{enumerate}

\[
((a \mathbf {M}) \cap (\mathfrak {b} \mathbf {M})) / (a \cap \mathfrak {b}) \mathbf {M}
\]

is pseudo-zero; give an example where it is not zero.

\begin{enumerate}
\def\labelenumi{\arabic{enumi}.}
\setcounter{enumi}{5}
\tightlist
\item
  Let \(k\) be a field, \(\mathbf{B}\) the polynomial ring \(k[\mathbf{X},\mathbf{Y}]\) and \(\mathbf{A}\) the submodule \(k[\mathbf{X}^2,\mathbf{XY},\mathbf{Y}^2]\) of \(\mathbf{B}\) .
\end{enumerate}

\begin{enumerate}
\def\labelenumi{(\alph{enumi})}
\item
  Show that A is an integrally closed Noetherian domain and that B is a finitely generated A-module.
\item
  Show that the ideal \(p = AX^{2} + AXY\) of A is a prime ideal of height 1 (and hence divisorial) but that the B-module \(p \otimes_{A} B\) is a torsion module \(\neq 0\) and hence not reflexive; the ideal pB of B is not divisorial, the canonical mapping \(p \otimes_{A} B \to pB\) is not injective and p is not a flat A-module.
\end{enumerate}

\begin{enumerate}
\def\labelenumi{\arabic{enumi}.}
\setcounter{enumi}{6}
\tightlist
\item
  Let A be an integrally closed Noetherian domain, E a finitely generated torsion-free A-module and \(\mathbf{E}^*\) its dual.
\end{enumerate}

\begin{enumerate}
\def\labelenumi{(\alph{enumi})}
\item
  Show that the canonical homomorphism \(\mathbf{E}^{*}\otimes_{\mathbf{A}}\mathbf{E}\to \mathrm{End}_{\mathbf{A}}(\mathbf{E})\) is a pseudo-isomorphism.
\item
  Deduce from (a) that for E to be a projective A-module, it is necessary and sufficient that \(E^{*} \otimes_{A} E\) be a reflexive A-module. (Note that, if \(E^{*} \otimes_{A} E\) is reflexive, the canonical homomorphism \(E^{*} \otimes_{A} E \to \operatorname{End}_{A}(E)\) is bijective.)
\end{enumerate}

¶* 8. Let A be an integrally closed Noetherian domain and \(\mathbf{M}_1, \mathbf{M}_2\) two finitely generated A-modules.

\begin{enumerate}
\def\labelenumi{(\alph{enumi})}
\item
  Show that the A-modules \(\operatorname{Tor}_i^{\mathbf{A}}(\mathbf{M}_1, \mathbf{M}_2)\) , \(\operatorname{Ext}_{\mathbf{A}}^{i}(\mathbf{M}_1, \mathbf{M}_2)\) are pseudo-zero for \(i \geqslant 2\) (reduce it to the case where \(\mathbf{A}\) is a principal ideal domain).
\item
  If \(\mathbf{M}_1\) is torsion-free, show that \(\operatorname{Tor}_1^{\mathrm{A}}(\mathbf{M}_1, \mathbf{M}_2)\) and \(\operatorname{Ext}_A^1(\mathbf{M}_1, \mathbf{M}_2)\) are pseudo-zero (same method).
\item
  If \(\mathbf{M}_1\) is a torsion A-module, show that (in the notation of no. 5):
\end{enumerate}

\[
\chi (\mathbf {M} _ {1} \otimes_ {\mathbf {A}} \mathbf {M} _ {2}) - \chi (\operatorname{Tor} _ {1} ^ {\mathbf {A}} (\mathbf {M} _ {1}, \mathbf {M} _ {2})) = r (\mathbf {M} _ {2}) \chi (\mathbf {M} _ {1})
\]

\[
\chi (\operatorname{Hom} _ {\mathbf {A}} (\mathbf {M} _ {1}, \mathbf {M} _ {2})) - \chi (\operatorname{Ext} _ {\mathbf {A}} ^ {1} (\mathbf {M} _ {1}, \mathbf {M} _ {2})) = - r (\mathbf {M} _ {2}) \chi (\mathbf {M} _ {1})
\]

\[
\chi (\mathrm{Hom} _ {\mathrm{A}} (\mathrm{M} _ {2}, \mathrm{M} _ {1})) - \chi (\mathrm{Ext} _ {\mathrm{A}} ^ {1} (\mathrm{M} _ {2}, \mathrm{M} _ {1})) = r (\mathrm{M} _ {2}) \chi (\mathrm{M} _ {1}).
\]

\begin{enumerate}
\def\labelenumi{(\alph{enumi})}
\setcounter{enumi}{3}
\tightlist
\item
  Show that for every ordered pair of finitely generated A-modules \(M_{1}\) , \(M_{2}\) :
\end{enumerate}

\[
\begin{array}{r} c (\mathbf {M _ {1}} \otimes_ {\mathbf {A}} \mathbf {M _ {2}}) - c (\mathrm{Tor} _ {\mathbf {1}} ^ {\mathbf {A}} (\mathbf {M _ {1}}, \mathbf {M _ {2}})) = r (\mathbf {M _ {1}}) c (\mathbf {M _ {2}}) + r (\mathbf {M _ {2}}) c (\mathbf {M _ {1}}) \\ c (\mathrm{Hom} _ {\mathbf {A}} (\mathbf {M _ {1}}, \mathbf {M _ {2}})) - c (\mathrm{Ext} _ {\mathbf {A}} ^ {1} (\mathbf {M _ {1}}, \mathbf {M _ {2}})) = r (\mathbf {M _ {1}}) c (\mathbf {M _ {2}}) - r (\mathbf {M _ {2}}) c (\mathbf {M _ {1}}). \end{array}
\]

\begin{enumerate}
\def\labelenumi{\arabic{enumi}.}
\setcounter{enumi}{8}
\tightlist
\item
  Let k be a field and A the polynomial ring \(k[X,Y]\) . On the A-module \(M = A^{2}\) consider the linear form f such that \(f(e_{1}) = X\) , \(f(e_{2}) = Y((e_{1}, e_{2})\) being the canonical basis). Show that the kernel L of f is a monogenous free A-module, but that the quotient M/L (which is isomorphic to an ideal of A) is not reflexive.
\end{enumerate}

¶ 10. Let A be a commutative ring. For every submodule R of a finitely generated free A-module \(L = A^n\) , let \(c_1(R)\) denote the ideal generated by the \(\langle x, x^* \rangle\) , where \(x\) runs through R and \(x^*\) runs through the dual \(L^*\) ; write \(c_k(R) = c_1\left(\operatorname{Im}\left(\bigwedge^k R\right)\right)\) , \(\operatorname{Im}\left(\bigwedge^k R\right)\) being the canonical image of the \(k\) -th exterior power of R in \(\bigwedge^k L\) . If \(R_1 \subset R_2\) are two submodules of L, then \(c_k(R_1) \subset c_k(R_2)\) for all \(k\) .

\begin{enumerate}
\def\labelenumi{(\alph{enumi})}
\item
  Let M be a finitely generated A-module; M is isomorphic to a quotient module L/R, where \(L = A^{n}\) for a suitable n.~Show that the sequence of ideals \((\mathfrak{a}_{k})_{k \geqslant 0}\) such that \(\mathfrak{a}_{k} = \mathfrak{c}_{n-k}(R)\) for \(k \leqslant n\) and \(\mathfrak{a}_{k} = A\) for k \textgreater{} n is independent of the expression of M in the form L/R. (Consider first the case where L/R and L/R' are isomorphic; then note that \(A^{n}/R\) is isomorphic to \(A^{n+h}/(R \times A^{h})\) for all h \textgreater{} 0.) We write \(\mathfrak{d}_{k}(M) = \mathfrak{a}_{k}\) for all \(k \geqslant 0\) and say that these are the determinantal ideals associated with M. Then \(\mathfrak{d}_{k}(M) \subset \mathfrak{d}_{k+1}(M)\) for \(k \geqslant 0\) .
\item
  If \(\mathbf{A}\) is a principal ideal domain and \(\mathbf{M} = \mathbf{L} / \mathbf{R}\) where \(\mathbf{L}\) is free and finitely generated, show that the ideals \(\mathfrak{c}_{k + 1}(\mathbb{R})(\mathfrak{c}_k(\mathbb{R}))^{-1}\) are the invariant factors of \(\mathbf{R}\) in \(\mathbf{L}\) .
\item
  Let \(r = \gamma(\mathbf{M})\) be the least of the cardinals of the systems of generators of \(\mathbf{M}\) and let \(r_0\) be the least integer \(h\) such that \(\mathfrak{d}_h(\mathbf{M}) = \mathbf{A}\) . Show that \(r_0 \leqslant r\) and give an example where \(r_0 < r\) (take \(\mathbf{A}\) to be a Dedekind domain).
\item
  If \(a\) is the annihilator of \(\mathbf{M}\) , show that \(\mathfrak{d}_0(\mathbf{M}) \subset \mathfrak{a}\) and \(\mathfrak{a}^{r - k} \subset \mathfrak{d}_k(\mathbf{M})\) for \(k \leqslant r = \gamma (\mathbf{M})\) .
\item
  Suppose that \(\mathbf{M}\) is a direct sum of submodules \(\mathbf{M}_i\) ( \(1 \leqslant i \leqslant h\) ). Show that
\end{enumerate}

\[
\mathfrak {d} _ {k} (\mathbf {M}) = \sum \mathfrak {d} _ {k _ {1}} (\mathbf {M} _ {1}) \dots \mathfrak {d} _ {k _ {n}} (\mathbf {M} _ {n})
\]

where the sum is over the finite sequences \((k_{i})_{1\leqslant i\leqslant h}\) such that \(\sum_{i = 1}^{h}k_{i} = k.\)

\begin{enumerate}
\def\labelenumi{(\alph{enumi})}
\setcounter{enumi}{5}
\tightlist
\item
  If N is a finitely generated submodule of M, then \(\mathfrak{d}_{k}(\mathbf{M}/\mathbf{N}) \subset \mathfrak{d}_{k}(\mathbf{M})\) for all \(k \geqslant 0\) . Show that:
\end{enumerate}

\[
\sum_ {j = 0} ^ {k} \mathfrak {d} _ {j} (\mathrm{N}) \mathfrak {d} _ {k - j} (\mathrm{M} / \mathrm{N}) \subset \mathfrak {d} _ {k} (\mathrm{M}).
\]

\begin{enumerate}
\def\labelenumi{(\alph{enumi})}
\setcounter{enumi}{6}
\tightlist
\item
  Let K be a field, A the polynomial ring K{[}X, Y, Z{]}, m the maximal ideal AX + AY + AZ and q the ideal generated by X, Y², YZ and Z². Consider the A-module M = A/q and its submodule N = m/q. Show that \(\mathfrak{d}_{0}(M) = \mathfrak{q}, \mathfrak{d}_{1}(M) = A, \mathfrak{d}_{0}(N) = m^{2}\) and \(\mathfrak{d}_{1}(N) = m\) .
\end{enumerate}

\begin{enumerate}
\def\labelenumi{\arabic{enumi}.}
\setcounter{enumi}{10}
\tightlist
\item
  Let A be a Krull domain and M, N two finitely generated lattices in a vector space V of finite rank n over the field of fractions of A. For all \(c \neq 0\) in A such that \(cN \subset M\) , consider the determinantal ideals \(\mathfrak{d}_{k}(M/cN)\) (Exercise 10) and write:
\end{enumerate}

\[
d _ {k} (\mathbf {N}, \mathbf {M}) = \operatorname{div} \left(\mathfrak {d} _ {k} (\mathbf {M} / c \mathbf {N})\right) - (n - k) \operatorname{div} (c).
\]

\begin{enumerate}
\def\labelenumi{(\alph{enumi})}
\item
  Show that \(d_{k}(\mathbf{N}, \mathbf{M})\) does not depend on the choice of c such that \(cN \subset M\) ; \(d_{k}(\mathbf{N}, \mathbf{M})\) is called the determinantal divisor of index k of N with respect to M.
\item
  \(d_{k}(\mathbf{N},\mathbf{M})\geqslant d_{k + 1}(\mathbf{N},\mathbf{M})\) and \(d_{n}(\mathbf{N},\mathbf{M}) = 0\) . Show that, if we write
\end{enumerate}

\[
e _ {k} (\mathrm{N}, \mathrm{M}) = d _ {n - k} (\mathrm{N}, \mathrm{M}) - d _ {n - k + 1} (\mathrm{N}, \mathrm{M}),
\]

then

\[
e _ {k} (\mathbf {N}, \mathbf {M}) \leqslant e _ {k + 1} (\mathbf {N}, \mathbf {M})
\]

for \(1 \leqslant k \leqslant n\) ; the divisors \(e_k(\mathbf{N}, \mathbf{M})\) are called the invariant factors of \(\mathbf{N}\) with respect to \(\mathbf{M}\) (reduce it to the case where \(\mathbf{A}\) is a principal ideal domain).

\begin{enumerate}
\def\labelenumi{(\alph{enumi})}
\setcounter{enumi}{2}
\item
  Show that \(d_0(\mathbf{N}, \mathbf{M}) = \chi(\mathbf{M}, \mathbf{N})\) (no. 5).
\item
  If M is a lattice in V but N is a lattice in a subspace W of V of rank q \textless{} n, the invariant factors of N with respect to \(M \cap W\) are called the invariant factors of N with respect to M. Show how Exercises 8 to 10 and 14 to 16 of Algebra, Chapter VII, §4 may be extended (assuming if need be in certain cases that A is an integrally closed Noetherian domain).
\end{enumerate}

\begin{enumerate}
\def\labelenumi{\arabic{enumi}.}
\setcounter{enumi}{11}
\tightlist
\item
  Let A be an integrally closed Noetherian domain and M, N two torsion-free finitely generated A-modules of the same rank r such that \(N \subset M\) ; let \(j: N \to M\) be the canonical injection.
\end{enumerate}

\begin{enumerate}
\def\labelenumi{(\alph{enumi})}
\tightlist
\item
  Show that, for all \(k\) , \(\bigwedge^{k}j\colon\bigwedge^{k}\mathrm{N}\to\bigwedge^{k}\mathrm{M}\) is pseudo-injective, that \(\operatorname{Coker}\left(\bigwedge^{k}j\right)\) is a torsion A-module and that:
\end{enumerate}

\[
\chi \left(\operatorname{Coker} \left(\bigwedge^ {k} j\right)\right) = \binom {r - 1} {k - 1} \chi (\operatorname{Coker} j).
\]

\begin{enumerate}
\def\labelenumi{(\alph{enumi})}
\setcounter{enumi}{1}
\tightlist
\item
  Using (a), show that, if M is a torsion-free finitely generated A-module, the torsion submodule of \(\bigwedge^{k}\mathbf{M}\) is pseudo-zero for all \(k\) and
\end{enumerate}

\[
c \left(\bigwedge^ {k} \mathbf {M}\right) = \binom {r - 1} {k - 1} c (\mathbf {M}),
\]

where \(\mathbf{M}\) is of rank \(r\) ; in particular, \(c\left(\bigwedge^{r}\mathbf{M}\right) = c(\mathbf{M})\) .

\begin{enumerate}
\def\labelenumi{\arabic{enumi}.}
\setcounter{enumi}{12}
\tightlist
\item
  Let \(k\) be a field and \(A\) the polynomial ring \(k[X, Y]\) .
\end{enumerate}

\begin{enumerate}
\def\labelenumi{(\alph{enumi})}
\item
  Let \(m\) be the maximal ideal \(AX + AY\) of \(A\) ; show that there exists a pseudo-isomorphism of \(m\) to \(A\) , but that there exists no pseudo-isomorphism of \(A\) to \(m\) .
\item
  Let \(\mathfrak{m}'\) be the maximal ideal \(\mathbf{A}(\mathbf{X} - 1) + \mathbf{A}\mathbf{Y}\) of \(\mathbf{A}\) ; show that \(c(\mathfrak{m}) = c(\mathfrak{m}') = 0\) , but that there exists no pseudo-isomorphism of \(\mathfrak{m}\) to \(\mathfrak{m}'\) nor of \(\mathfrak{m}'\) to \(\mathfrak{m}\) .
\item
  Let \(\mathfrak{p} = \mathbf{AX}\) , \(\mathfrak{q} = \mathbf{AY}\) , which are prime ideals of height 1. In the A-module \(\mathbf{L} = \mathbf{A}^2\) , consider the submodules \(\mathbf{M} = \mathfrak{pe}_1 \oplus \mathfrak{qe}_2\) , \(\mathbf{N} = \mathbf{Ae}_1 \oplus \mathfrak{pqe}_2\) ( \(e_1, e_2\) being the vectors of the canonical basis of \(\mathbf{L}\) ). Show that \(\mathbf{M}\) and \(\mathbf{N}\) are isomorphic and that there exist a pseudo-isomorphism from \(\mathbf{L} / \mathbf{M}\) to \(\mathbf{L} / \mathbf{N}\) and a pseudo-isomorphism from \(\mathbf{L} / \mathbf{N}\) to \(\mathbf{L} / \mathbf{M}\) but that there exists no pseudo-isomorphism from \(\mathbf{L}\) to itself mapping \(\mathbf{M}\) to \(\mathbf{N}\) or \(\mathbf{N}\) to \(\mathbf{M}\) (observe that a pseudo-isomorphism from \(\mathbf{L}\) to itself is necessarily an automorphism of \(\mathbf{L}\) , with the aid of Proposition 10).
\end{enumerate}

¶ 14. Let A be a Prüferian domain (§ 2, Exercise 12) and M, N two finitely generated lattices in a vector space V of finite rank n over the field of fractions of A. For all \(c \neq 0\) in A such that \(cN \subset M\) , consider the determinantal ideals \(\mathfrak{d}_{k}(M/cN)\) (Exercise 10) and write:

\[
\mathfrak {d} _ {k} (\mathbf {N}, \mathbf {M}) = c ^ {k - n} \mathfrak {d} _ {k} (\mathbf {M} / c \mathbf {N}).
\]

\begin{enumerate}
\def\labelenumi{(\alph{enumi})}
\item
  Show that \(\mathfrak{d}_k(\mathbf{N},\mathbf{M})\) does not depend on the choice of \(c\) such that \(c\mathbf{N}\subset \mathbf{M}\) ; these are called the determinantal ideals of \(\mathbf{N}\) with respect to \(\mathbf{M}\) .
\item
  \(\mathfrak{d}_k(\mathbf{N},\mathbf{M})\subset \mathfrak{d}_{k + 1}(\mathbf{N},\mathbf{M})\) and \(\mathfrak{d}_n(\mathbf{N},\mathbf{M}) = \mathbf{A}\) . Show that, if we write \(\mathfrak{e}_k(\mathbf{N},\mathbf{M}) = \mathfrak{d}_{n - k}(\mathbf{N},\mathbf{M})(\mathfrak{d}_{n - k + 1}(\mathbf{N},\mathbf{M}))^{-1}\) , the integral ideals \(\mathfrak{e}_k(\mathbf{N},\mathbf{M})\) are such that \(\mathfrak{e}_k(\mathbf{N},\mathbf{M})\supset \mathfrak{e}_{k + 1}(\mathbf{N},\mathbf{M})\) for \(1\leqslant k\leqslant n\) ; the finitely generated ideals \(\mathfrak{e}_k(\mathbf{N},\mathbf{M})\) are called the invariant factors of \(\mathbf{N}\) with respect to \(\mathbf{M}\) (consider the \(\mathrm{A_m}\) -modules \(\mathrm{M_m}\) and \(\mathrm{N_m}\) for every maximal ideal \(\mathfrak{m}\) of \(\mathrm{A}\) ).
\item
  If M is a lattice of V but N is a lattice in a subspace W of V of rank q \textless{} n, the invariant factors of N with respect to M ∩ W are called the invariant factors of N with respect to M. Show how Exercises 8 to 10 and 14 to 16 of Algebra, Chapter VII, § 4 may be extended.
\end{enumerate}

\begin{enumerate}
\def\labelenumi{\arabic{enumi}.}
\setcounter{enumi}{14}
\tightlist
\item
  Let A be the polynomial ring Z{[}X, Y, T, U, V, W{]}, L the free A-module \(A^{4}\) , \((e_{i})_{1 \leq i \leq 4}\) its canonical basis and M the submodule of L generated by the four vectors \(Xe_{1}\) , \(Ye_{2}\) , \(Te_{3} + Ue_{4}\) , \(Ve_{3} + We_{4}\) . Show that
\end{enumerate}

\[
\mathfrak {d} _ {1} (\mathrm{L} / \mathrm{M}) \mathfrak {d} _ {3} (\mathrm{L} / \mathrm{M}) \notin (\mathfrak {d} _ {2} (\mathrm{L} / \mathrm{M})) ^ {2}
\]

(compare with Exercise 14 (b)).

¶ 16. (a) Let A be a Prüferian domain and M a finitely generated lattice in a vector space V of finite rank n over the field of fractions K of A. Show that there exists a basis \((e_{i})_{1 \leqslant i \leqslant n}\) of V and n fractional ideals \(a_{i}\) ( \(1 \leqslant i \leqslant n\) ) such that M is equal to the direct sum of the \(a_{i}e_{i}\) . (Argue by induction on n: if \((u_{i})_{1\leq i\leq n}\) is a basis of V, consider the finitely generated fractional ideal \(b_{1}\) generated by the \(u_{1}\) -coordinates of the elements of M; by considering the A-module \(b_{1}^{-1}M\) , reduce it the case where \(b_{1}=A\) .)

Deduce that, if W is a subspace of V, M ∩ W is a direct factor of M.

\begin{enumerate}
\def\labelenumi{(\alph{enumi})}
\setcounter{enumi}{1}
\tightlist
\item
  Suppose that M is a submodule of \(A^{n}\) equal to the direct sum of the \(a_{i}e_{i}\) , where the \(a_{i}\) are finitely generated (integral) ideals of A and \((e_{i})_{1\leqslant i\leqslant n}\) the canonical basis of \(A^{n}\) . Show that there exists a basis \((u_{i})_{1\leqslant i\leqslant n}\) of \(A^{n}\) and ideals \(b_{i}\) of A such that M is the direct sum of the \(b_{i}u_{i}\) and \(b_{i}\) divides \(b_{i+1}\) for \(1\leqslant i\leqslant n-1\) . (Reduce it to the case where n=2 and \(a_{1}+a_{2}=A\) .)
\end{enumerate}

\begin{enumerate}
\def\labelenumi{\arabic{enumi}.}
\setcounter{enumi}{16}
\tightlist
\item
  Let k be a field, A the polynomial ring \(k[X, Y]\) , L the free A-module \(A^{2}\) and \((e_{1}, e_{2})\) the canonical basis of L. Let M be the submodule of L generated by the two vectors \((X + 1)e_{1} + Ye_{2}, Ye_{1} + Xe_{2}\) . Show that there exists no pseudo-isomorphism from L to itself, whose restriction to M is a pseudo-isomorphism from M to the submodule N of the form \(au + bv\) where \((u, v)\) is a basis of the A-module L and a, b two ideals of A, or whose restriction to N is a pseudo-isomorphism from N to M. (By considering the determinantal ideals (Exercise 10) and noting that M is reflexive, attention may be confined to the case where N would also be reflexive and then necessarily a = A, b = AP, where
\end{enumerate}

\[
\mathrm{P} (\mathrm{X}, \mathrm{Y}) = \mathrm{X} (\mathrm{X} + 1) - \mathrm{Y} ^ {2}
\]

is an extremal element of A; show finally that, for every basis \((u, v)\) of L, M ∩ Au can neither contain Au nor be contained in APu.)

¶ 18. Let A be a Dedekind domain, K its field of fractions and M, N two lattices in a vector space V of finite rank n over K. Let \(\mathfrak{e}_{k} = \mathfrak{e}_{k}(N, M)\) be the invariant factors of N with respect to M (Exercise 14). Show that there exists a basis \((u_{i})_{1 \leqslant i \leqslant n}\) of V such that M is equal to a direct sum \(\bigoplus_{i} a_{i} u_{i}\) , where the \(a_{i}\) are fractional ideals, and N is equal to the direct sum \(\bigoplus_{i} e_{i} a_{i} u_{i}\) . (Use the theory of finitely generated modules over a principal ideal domain (Algebra, Chapter VII, § 4, no. 2, Theorem 1) and the approximation theorem in the unimodular group \(\mathbf{SL}(n, \mathbf{A})\) (§ 2, no. 4).)

Extend this result to the case where A is the union of a right directed family of subrings which are Dedekind domains (for example the integral closure of a Dedekind domain in the algebraic closure of its field of fractions).

\begin{enumerate}
\def\labelenumi{\arabic{enumi}.}
\setcounter{enumi}{18}
\item
  Let A be a Dedekind domain and a, b fractional ideals of A. Show that the A-modules A ⊕ ab and a ⊕ b are isomorphic.
\item
  Let A be a Dedekind domain, P a projective A-module and N a submodule of P. Show that N is projective and a direct sum of modules isomorphic to ideals of A. (Attention may be confined to the case where P is free. Then proceed as in Exercise 16, using a transfinite induction argument based on that of Algebra, Chapter VII, § 3, Theorem 1.)
\item
  Let P be a projective module over a Dedekind domain A. Show that, if P is not finitely generated, it is free. (By virtue of Exercise 20, we may write P as an infinite direct sum \(\bigoplus_{\lambda} (\mathfrak{b}_{\lambda} \oplus \mathfrak{c}_{\lambda})\) , where, for each \(\lambda, \mathfrak{b}_{\lambda}\) and \(\mathfrak{c}_{\lambda}\) are ideals of A. Using Exercise 19, \(\mathbf{P} = \mathbf{L} \oplus \mathbf{Q}\) , where \(\mathbf{L}\) is isomorphic to \(\mathbf{A}^{(I)}\) (I infinite) and \(\mathbf{Q} = \bigoplus_{\alpha \in I} \mathfrak{a}_{\alpha}\) , where the \(\mathfrak{a}_{\alpha}\) are ideals of A. Then apply Exercise 3 of Algebra, Chapter II, § 2.)
\end{enumerate}

¶ 22. Let A be a local ring, m its maximal ideal and E a finitely generated A-module.

\begin{enumerate}
\def\labelenumi{(\alph{enumi})}
\item
  If \(r = \gamma(\mathrm{E})\) is the least of the cardinals of a system of generators of \(\mathbf{E}\) , \(r\) is also the least integer \(h\) such that the determinantal ideal \(\mathfrak{d}_h(\mathbf{E})\) is equal to \(\mathbf{A}\) and also the least of the integers \(k\) such that \(\bigwedge^{k+1} \mathbf{E} = 0\) (note that \(r\) is the rank of \(\mathbf{E}/\mathfrak{m}\mathbf{E}\) over \(\mathbf{A}/\mathfrak{m}\) ).
\item
  If we write \(e(E) = \mathfrak{d}_{r-1}(E)\) , show that \(\bigwedge^{r} E\) is isomorphic to \(A / e(E)\) . For an ideal \(a\) of \(A\) to contain \(e(E)\) , it is necessary and sufficient that \(E / aE\) be a free \((A / a)\) -module (reduce it to the case where \(a = 0\) ).
\item
  If \(\mathfrak{e}(\mathbf{E})\) is a principal ideal \(A\alpha\) , show that E contains a direct factor isomorphic to \(A/A\alpha\) (write E as a quotient of \(A^{r}\) by a submodule R and note that there is in R an element of the form \(\alpha y\) , where y is an element of a basis of \(A^{r}\) ).
\item
  Let \(\mathfrak{b}_h^{\prime}(\mathbf{E})\) be the annihilator of \(\bigwedge^{h+1}\mathbf{E}\) for \(0 \leqslant h \leqslant r - 1\) . Show that the following conditions are equivalent:
\end{enumerate}

\((\alpha)\) The ideals \(\mathfrak{d}_h(\mathbf{E})\) ( \(0 \leqslant h \leqslant r - 1\) ) are principal.

(β) The ideals \(\mathfrak{d}_h^\prime (\mathrm{E})\) \((0\leqslant h\leqslant r - 1)\) are principal.

\((\gamma)\) E is the direct sum of r modules isomorphic to A/A \(\lambda_{h}\) , where \(\lambda_{h+1}\) divides \(\lambda_{h}\) for \(0 \leqslant h \leqslant r - 1\) .

Moreover, if these conditions hold, then \(\mathfrak{d}_h'(\mathrm{E}) = \mathrm{A}\lambda_h\) for \(0 \leqslant h \leqslant r - 1\) and \(\mathfrak{d}_h(\mathrm{E}) = \lambda_h\lambda_{h+1} \ldots \lambda_{r-1}\mathrm{A}\) for \(0 \leqslant h \leqslant r - 1\) . (Proceed by induction on \(r\) , using (c).)

\begin{enumerate}
\def\labelenumi{(\alph{enumi})}
\setcounter{enumi}{4}
\tightlist
\item
  Suppose further that A is an integral domain. Show that, if p is a prime ideal of A such that E/pE is a free (A/p)-module and \(E_{p}\) a free \(A_{p}\) -module, then E is a free A-module (using (b), show that \(\gamma(E_{p}) = \gamma(E)\) and that \(e(E_{p}) = e(E)\) ).
\end{enumerate}

\begin{enumerate}
\def\labelenumi{\arabic{enumi}.}
\setcounter{enumi}{22}
\item
  Let A be a ring, E a finitely generated A-module, r the least of the integers k such that \(\bigwedge^{k+1}E=0\) and \(e(E)\) the annihilator of \(\bigwedge^{r}E\) . Show that, for every ideal \(a\supset e(E)\) in A, the \((A/a)\) -module E/aE is flat (reduce it to the case where A is a local ring and use Exercise 22 (b)). Deduce that, if r is the Jacobson radical of A and if, for every maximal ideal m of A, the rank of the vector \((A/m)\) -space E/mE is the same, then E/rE is a flat \((A/r)\) -module.
\item
  Let A be an integrally closed Noetherian domain. For a lattice M (with respect to A) to be reflexive, it is necessary and sufficient that it satisfy the following condition: for every ordered pair \((a, b)\) of elements of A such that \(a \neq 0\) and the homothety of ratio b on A/aA is injective, the homothety of ratio b on M/aM is then injective. (To see that the condition is necessary, consider two elements x, y of M such that \(ax + by = 0\) and show that, for all \(p \in P(A)\) , \(y \in aM_{p}\) . To show that the condition is sufficient, use criterion (c) of Theorem 2 and observe (using Proposition 8 of §1, no. 4) that the hypothesis may be written as
\end{enumerate}

\[
\operatorname{Ass} (\mathbf {M} / a \mathbf {M}) = \operatorname{Ass} (\mathbf {A} / a \mathbf {A})
\]

for all \(a \neq 0\) in \(\mathbf{A}\) and that for every module \(\mathbf{E}\) such that \(\mathbf{M} \subset \mathbf{E} \subset \mathbf{V}\) there exists \(a \neq 0\) in \(\mathbf{A}\) such that \(a\mathbf{M} \subset a\mathbf{E} \subset \mathbf{M}\) .

\begin{enumerate}
\def\labelenumi{\arabic{enumi}.}
\setcounter{enumi}{24}
\tightlist
\item
  Let \(A = \bigoplus_{n \geq 0} A_n\) be a Noetherian graded ring with positive degrees. Show that, if A is a factorial domain, then \(A_0\) is a factorial domain and the \(A_n\) are reflexive. (To show that \(A_0\) is a factorial domain, use criterion (c) of §3, no. 2, Theorem 1; to show that the \(A_n\) are reflexive, use Exercise 24.)
\end{enumerate}

¶ 26. Let A be a Noetherian ring and M a finitely generated A-module. For the symmetric algebra S(M) to be a factorial domain, it is necessary and sufficient that A be a factorial domain and that the \(S^{n}(M)\) be reflexive A-modules. (To see that the condition is sufficient, observe first that, if T = A - \{0\}, S(M) is identified with a subring of \(T^{-1}S(M)\) ; then show that every extremal element p of A is extremal in S(M), by reducing it to proving that, if p divides a product xy of two homogeneous elements in S(M), it divides one of them; finally, apply Proposition 3 of § 3, no. 4.)

